\section{Data}

\paragraph{Source, scope and what the label means.}
The data are 30,000 existing cardholders of one Taiwanese bank (\emph{Default of Credit Card Clients}, UCI Machine Learning Repository, \url{https://archive.ics.uci.edu/dataset/350/default+of+credit+card+clients}, doi:10.24432/C55S3H, CC~BY~4.0 [18]; described in [19]). We predict default on the October 2005 payment from April--September 2005 histories, so a flag triggers an account-management action on a current client (a call, a reminder, a limit review), not an approval. The monthly variables precede the label; the limit and demographics are undated snapshots, assumed recorded by September 2005 (if the bank cut limits after missed payments, the limit would partly record the outcome). Removing 35 rows that repeat another client's record exactly, label included (mostly dormant profiles), leaves 29,965 clients, 29 predictors (Table~\ref{tab:a-variables}) and 22.1\% defaulters: ``never default'' is 77.9\% accurate yet catches nobody. Code: \href{https://github.com/jingwen013-design/mh6805-credit-default}{\texttt{github.com/jingwen013-design/mh6805-credit-default}} (\texttt{code/run\_all.sh} regenerates every number).

\paragraph{Repayment codes are states, not a scale.}
PAY\_1--PAY\_6 record each month's status: $-2$ no consumption, $-1$ paid in full, 0 revolving credit, $k \ge 1$ = $k$ months late (UCI documents only $-1$ and 1--9; the readings of $-2$ and 0 are widely used, not documented). Non-delinquent codes have no order: September default is 13.7\% at $-2$, 16.9\% at $-1$ and 12.9\% at 0 (training rows). Code 1 also changes meaning: it appears 2,540 times in September but 21 times in the five earlier months, mixing dormant accounts with clients recovering from two months late (Table~\ref{tab:a-paycodes-rates}), so compare codes within a month. A scorecard should weight these states separately (\S4).

\paragraph{Where the risk gradients are.}
Default rates change steeply with recent arrears and the limit, not with bill size. September default climbs from 12.9\% at revolving credit to 67.6\% at two months late, two payments already missed (code 1 is the mixture above). It falls steadily as the limit rises, but the bank sets the limit: a high limit marks a client it already judged safe, so raising a limit would not make anyone safer. Bill size alone ranks no better than chance (univariate area under the ROC curve, AUC, 0.50--0.52). Utilisation (bill over limit) is U-shaped: default is lowest for light users and highest near the limit, where no buffer is left. It is also above average with no positive balance, but most such clients carry the odd September code 1, so that arm is of uncertain meaning. Sex and age move the rate by at most 8 points (training rows; Table~\ref{tab:a-gradients}, Figures~\ref{fig:eda-default-rate-by-pay1}--\ref{fig:eda-default-rate-util-delay}).

\paragraph{What happens before the split.}
Three fixed steps estimate nothing, so they precede the split; everything estimated (scaling, one-hot levels, resampling, thresholds) is re-fitted inside each fold's \texttt{Pipeline} (Table~\ref{tab:a-quality}). (i)~An a-priori mapping of EDUCATION and MARRIAGE codes the UCI dictionary does not list (\{0,5,6\} and 0 to ``others''), so the ``others'' education bucket (343 training clients, 6.7\% default) is a coding artefact. (ii)~Six engineered features: two utilisation ratios, a payment ratio capped at 5, months in arrears, worst delay and balance trend (the delay counters count $-2$, $-1$ and 0 as 0). (iii)~Removal of the 35 repeated rows. Filling the payment ratio with 0 when all six bills are non-positive makes owing nothing look like paying nothing. The fill is not neutral: these 633 training clients default at 37.0\% against 21.7\%, so a bank must document it.
